% Source PDF pages 114--117; manual pages 2-85--2-88.
\subsection{Intercepts and Proportional Navigation}\label{sec:intercepts-pn}

TAOS can compute intercept trajectories by optimization, predictive guidance, or
proportional navigation. Optimization constrains trajectory final conditions to meet
the target while varying the trajectory shape to minimize or maximize another
quantity such as time or velocity. Predictive guidance and proportional navigation
instead construct an intercept trajectory directly at each integration point.

Predictive guidance assumes a nonaccelerating target, estimates a straight-line
intercept point, and commands the vehicle toward it. Proportional navigation tries to
hold the line-of-sight direction constant by commanding acceleration proportional
to line-of-sight angular rate.

These are idealized trajectory methods. TAOS does not model the interceptor control
system or the sensors used to detect and track the target, so the results are not
intended to predict miss distance. They are useful for overall performance estimates
such as time and range to intercept.

\taossubhead{Predictive Guidance}

In the geometry of \cref{fig:predictive-intercept-geometry}, vehicle 1 is the
interceptor and vehicle 2 is the target. If both vehicles are treated as
nonaccelerating and occupy the same position at intercept time $t_i$, then
\begin{equation}
  \vec r_1(t)+(t_i-t)\vec V_1
  =\vec r_2(t)+(t_i-t)\vec V_2.
  \label{eq:predictive-intercept-position-equality}
\end{equation}
The time to intercept is estimated as
\begin{equation}
  t_i-t=\frac{\lVert\vec r_2-\vec r_1\rVert}
              {\lVert\vec V_1-\vec V_2\rVert},
  \label{eq:predictive-time-to-intercept}
\end{equation}
and the predicted intercept point is
\begin{equation}
  \vec r_i=\vec r_2(t)+\vec V_2
  \frac{\lVert\vec r_2-\vec r_1\rVert}
       {\lVert\vec V_1-\vec V_2\rVert}.
  \label{eq:predictive-intercept-point}
\end{equation}

\begin{figure}[htbp]
  \centering
  \begin{tikzpicture}[>=Latex,x=0.95cm,y=0.95cm,line cap=round]
  \coordinate (o) at (0,0);
  \coordinate (target) at (5.0,4.15);
  \coordinate (interceptor) at (5.35,1.05);
  \coordinate (impact) at (7.1,4.15);

  \draw[->,line width=0.75pt] (o) -- (target) node[midway,above left] {$\vec r_2(t)$};
  \draw[->,line width=0.75pt] (o) -- (interceptor) node[midway,below] {$\vec r_1(t)$};
  \draw[->,line width=0.75pt] (o) -- (impact) node[midway,above] {$\vec r_i$};
  \draw[->,line width=0.75pt] (target) -- (impact) node[midway,above] {$\vec V_2$};
  \draw[->,line width=0.75pt] (interceptor) -- ++(1.45,0.72) node[right] {$\vec V_1$};
  \draw[->,line width=0.75pt] (interceptor) -- (impact) node[midway,right] {$\Delta\vec r$};

  \fill (o) circle (2.1pt);
  \fill (target) circle (2.1pt);
  \fill (interceptor) circle (2.1pt);
  \fill (impact) circle (2.1pt);

  \draw[->] (o) -- ++(0.8,0);
  \draw[->] (o) -- ++(0,0.8);
  \draw[->] (o) -- ++(-0.45,-0.45);
  \node[left,align=right] at (-0.1,0.55) {ECFC\\system};
  \node[above left] at (target) {Target};
  \node[below right] at (interceptor) {Interceptor};
  \node[right,align=left] at (impact) {Intercept\\point};
\end{tikzpicture}

  \caption{Intercept geometry.}
  \label{fig:predictive-intercept-geometry}
\end{figure}

TAOS determines geodetic horizontal and vertical flight-path angles for the straight
line from the interceptor to $\vec r_i$ and uses them as guidance rules in the
solution process of \cref{sec:control-variable-solution}. The calculation is
implemented in \texttt{intercept} in the \texttt{derivs} module.

\taossubhead{Proportional Navigation}

Proportional navigation commands accelerations normal to the line of sight and
proportional to its rate. In \cref{fig:proportional-navigation-geometry}, vehicle 1
is the interceptor, vehicle 2 is the target, $\Delta\vec r$ is their relative position,
and $\Delta\vec V$ is their relative velocity.

\begin{figure}[htbp]
  \centering
  \begin{tikzpicture}[>=Latex,x=1.0cm,y=1.0cm,line cap=round]
  \coordinate (o) at (0,0);
  \coordinate (proj) at (6.3,0);
  \coordinate (target) at (6.3,4.25);
  \coordinate (yend) at (3.55,0.95);

  \draw[->] (o) -- ++(-0.8,-0.85) node[below left] {$\hat x_{\oplus}$};
  \draw[->] (o) -- (yend) node[above] {$\hat y_{\oplus}$};
  \draw[->] (o) -- ++(0,3.65) node[left] {$\hat z_{\oplus}$};
  \draw[->,line width=0.85pt] (o) -- (target) node[midway,above left] {$\Delta\vec r$};
  \draw[->] (o) -- (proj);
  \draw[densely dashed] (proj) -- (target);
  \draw[->] (target) -- ++(-1.45,0.2) node[above left] {$\Delta\vec V$};

  \draw[->] (1.25,0) arc[start angle=0,end angle=-34,radius=1.25];
  \node at (1.35,-0.55) {$\Lambda_y$};
  \draw[->] (4.75,0) arc[start angle=0,end angle=34,radius=1.55];
  \node at (4.55,0.95) {$\Lambda_p$};

  \fill (o) circle (2pt);
  \fill (target) circle (2pt);
  \node[left,align=right] at (-0.2,0.55) {Vehicle 1\\(interceptor)};
  \node[right,align=left] at (target) {Vehicle 2\\(target)};
  \node[right,align=left] at (proj) {$\Delta\vec r$ projected\\onto the $x$--$y$ plane};
\end{tikzpicture}

  \caption{Proportional-navigation geometry.}
  \label{fig:proportional-navigation-geometry}
\end{figure}

The line-of-sight yaw and pitch angles relative to ECFC coordinates are defined by
\begin{equation}
  \tan\Lambda_y
  =\frac{\Delta\vec r\cdot\hat y_{\oplus}}
         {\Delta\vec r\cdot\hat x_{\oplus}},
  \label{eq:line-of-sight-yaw}
\end{equation}
\begin{equation}
  \tan\Lambda_p
  =\frac{\Delta\vec r\cdot\hat z_{\oplus}}
  {\sqrt{(\Delta\vec r\cdot\hat x_{\oplus})^2
        +(\Delta\vec r\cdot\hat y_{\oplus})^2}}.
  \label{eq:line-of-sight-pitch}
\end{equation}
Differentiation gives
\begin{equation}
  \dot\Lambda_y=
  \frac{
  (\Delta\vec r\cdot\hat x_{\oplus})(\Delta\vec V\cdot\hat y_{\oplus})
  -(\Delta\vec r\cdot\hat y_{\oplus})(\Delta\vec V\cdot\hat x_{\oplus})}
  {(\Delta\vec r\cdot\hat x_{\oplus})^2
  +(\Delta\vec r\cdot\hat y_{\oplus})^2},
  \label{eq:line-of-sight-yaw-rate}
\end{equation}
\begin{equation}
\begin{aligned}
  \dot\Lambda_p={}&
  \Bigl[(\Delta\vec V\cdot\hat z_{\oplus})
   \bigl((\Delta\vec r\cdot\hat x_{\oplus})^2
        +(\Delta\vec r\cdot\hat y_{\oplus})^2\bigr)\\
  &\quad-(\Delta\vec r\cdot\hat z_{\oplus})
   \bigl((\Delta\vec r\cdot\hat x_{\oplus})
         (\Delta\vec V\cdot\hat x_{\oplus})
        +(\Delta\vec r\cdot\hat y_{\oplus})
         (\Delta\vec V\cdot\hat y_{\oplus})\bigr)\Bigr]\\
  &\Bigg/\Bigl[
  \sqrt{(\Delta\vec r\cdot\hat x_{\oplus})^2
        +(\Delta\vec r\cdot\hat y_{\oplus})^2}
  \bigl((\Delta\vec r\cdot\hat x_{\oplus})^2
       +(\Delta\vec r\cdot\hat y_{\oplus})^2
       +(\Delta\vec r\cdot\hat z_{\oplus})^2\bigr)\Bigr].
\end{aligned}
\label{eq:line-of-sight-pitch-rate}
\end{equation}
The desired normal accelerations are
\begin{equation}
  n_y=N'V_c\dot\Lambda_y,
  \label{eq:proportional-navigation-yaw-acceleration}
\end{equation}
\begin{equation}
  n_p=N'V_c\dot\Lambda_p,
  \label{eq:proportional-navigation-pitch-acceleration}
\end{equation}
where $N'$ is the navigation constant, typically between two and five, and
closure velocity is
\begin{equation}
  V_c=-\frac{\Delta\vec V\cdot\Delta\vec r}
             {\lVert\Delta\vec r\rVert}.
  \label{eq:proportional-navigation-closure-velocity}
\end{equation}

The source constructs a line-of-sight coordinate system with its $x$ axis aligned
with $\Delta\vec r$. The two printed rotation matrices are
\begin{equation}
  [R_p]=
  \begin{bmatrix}
    \cos\Lambda_p & \sin\Lambda_p & 0\\
    -\sin\Lambda_p & \cos\Lambda_p & 0\\
    0 & 0 & 1
  \end{bmatrix},
  \label{eq:source-los-pitch-rotation}
\end{equation}
\begin{equation}
  [R_y]=
  \begin{bmatrix}
    \cos\Lambda_y & 0 & \sin\Lambda_y\\
    0 & 1 & 0\\
    -\sin\Lambda_y & 0 & \cos\Lambda_y
  \end{bmatrix}.
  \label{eq:source-los-yaw-rotation}
\end{equation}
Their product, transposed for acceleration transformation, is
\begin{equation}
  [T_{\oplus\leftarrow\mathrm{LOS}}]=
  \begin{bmatrix}
    \cos\Lambda_p\cos\Lambda_y & -\sin\Lambda_p\cos\Lambda_y & -\sin\Lambda_y\\
    \sin\Lambda_p & \cos\Lambda_p & 0\\
    \cos\Lambda_p\sin\Lambda_y & -\sin\Lambda_p\sin\Lambda_y & \cos\Lambda_y
  \end{bmatrix},
  \label{eq:los-to-ecfc-acceleration-matrix}
\end{equation}
and the commanded ECFC acceleration components are
\begin{equation}
  \begin{bmatrix}
    \vec a_{\oplus}\cdot\hat x_{\oplus}\\
    \vec a_{\oplus}\cdot\hat y_{\oplus}\\
    \vec a_{\oplus}\cdot\hat z_{\oplus}
  \end{bmatrix}
  =[T_{\oplus\leftarrow\mathrm{LOS}}]
  \begin{bmatrix}0\\n_p\\n_y\end{bmatrix}.
  \label{eq:proportional-navigation-ecfc-acceleration}
\end{equation}

The component along the line-of-sight $x$ axis is set to zero because impact speed
is not the primary concern; the normal components turn the trajectory toward the
target. The resulting ECFC acceleration can be converted to desired horizontal and
vertical flight-path-angle rates with
\cref{eq:geodetic-flight-path-angle-rate,eq:geodetic-heading-rate}. The
Newton--Raphson guidance loop then solves for control values that produce those
rates. The proportional-navigation calculations are implemented by
\texttt{prop\_nav} in the \texttt{derivs} module.
